\documentclass[11pt]{article}
\usepackage[margin=0.9in]{geometry}
\usepackage{graphicx,amsmath,amssymb,booktabs,hyperref,xcolor,float}
\graphicspath{{./}}
\hypersetup{colorlinks=true,linkcolor=blue,urlcolor=blue}
\setlength{\parskip}{6pt}
\title{Implicit-function-theorem (IFT) Newton on the strict-$h{=}0$ K=3 REE \\
       in private/public log-odds coordinates}
\author{Fixed-Point Factory}
\date{June 2026}
\begin{document}
\maketitle

\begin{abstract}
We replace the damped-Picard/Anderson FP iteration on the strict-$h{=}0$
K=3 operator (in $(L_{\rm priv}, L_{\rm pub})$ coordinates, unbounded
cube, density-weighted residual) with an IFT-based Newton solver.
The implicit function theorem gives the Newton direction
$\delta P = -(\partial \Phi/\partial P - I)^{-1} (\Phi(P) - P)$.
We use \texttt{scipy.optimize.newton\_krylov} (LGMRES + Wolfe line
search; finite-diff Jacobian-vector products supply the implicit
$\partial \Phi/\partial P$ without forming it). We also tested
\texttt{scipy.optimize.root method='hybr'} (MINPACK HYBRD with explicit
numerical Jacobian + Broyden updates); HYBRD was either captured by
the trivial $P{=}0.5$ basin (logit-space variant) or wandered out of
the unit interval (P-space, even with bound penalties), so the
Newton-Krylov variant is the one reported.

\smallskip
\noindent\textbf{Setup.}
\begin{itemize}\itemsep1pt
  \item Grid: $G = 11$, $\xi$-uniform private log-odds, $\xi_{\max} = 0.85$
        ($u$ extends to $\pm 5\sigma$).
  \item $p$-grid: bootstrap symmetrized $L_{\rm pub}$ (frozen during Newton).
  \item Operator: strict-$h{=}0$ co-area + POU; correct K=3 CRRA clearing
        $\sum_k (R_k - 1)/((1-p) + R_k p) = 0$ with
        $R_k = \exp((\operatorname{logit}\mu_k - \operatorname{logit}p)/\gamma)$.
  \item Solver: \texttt{scipy.optimize.newton\_krylov} (LGMRES + Wolfe
        line search), $f\_{\rm tol} = 10^{-12}$, \texttt{maxiter} 30 Newton steps.
  \item Symmetry projection (permutation + sign-flip) applied to every $\Phi$ output.
  \item Warm-start: Picard FP at same $\tau$ from
        \texttt{logodds\_ladder/} (informative branch, not trivial $P{=}0.5$).
\end{itemize}

\smallskip
\noindent\textbf{Ladder.} 50 log-spaced $\tau \in [0.001, 1.2]$,
$\gamma = 100$, $G = 11$. Each cell warm-started from its corresponding
Picard FP.
\end{abstract}

\section{Why Newton from a cold start fails}

A naive Newton-via-IFT from the symmetric init $P = \sigma(\tau \sum u)$
converges to the \textbf{trivial $P \equiv 0.5$ FP} in $\sim 30$
$F$-evaluations, for every $\tau$ in the ladder. Reason: $P = 0.5$
exactly satisfies the K=3 CRRA clearing condition (all $\mu_k = 0.5$
gives $\sum_k (R_k-1)/((1-p)+R_k p) = 0$ at $p = 0.5$), and it's
inside the basin of attraction from the (small-amplitude) symmetric
init. Newton's local-quadratic convergence drags you into the nearest
FP regardless of economic interest.

The K=3 REE is multi-FP: trivial uninformative $P = 0.5$ and
informative branches. To stay on the informative branch we warm-start
Newton from the Picard FP found in the previous ladder run.

\section{Results}

\begin{figure}[H]\centering
  \includegraphics[width=\linewidth]{ladder_newton.png}
  \caption{Newton-via-IFT $\tau$-ladder, $\gamma = 100, G = 11$,
    warm-started from Picard FPs.
    \textbf{Left}: residual norms vs $\tau$ ($F_\infty$,
    density-weighted $F_w$, sign-flip $sym_{\rm err}$).
    \textbf{Middle}: equilibrium price range $[P_{\min}, P_{\max}]$.
    \textbf{Right}: Newton-via-IFT $F$-eval count per cell.}
\end{figure}

\section{Reading the result}

\begin{itemize}\itemsep1pt
  \item \textbf{Low $\tau$ ($\le 0.002$)}: trivial FP $P \equiv 0.5$
        wins (Picard warm-start was very close to it; Newton snaps in).
        $F_\infty$ at machine $\varepsilon$.
  \item \textbf{Transition $\tau \in [0.002, 0.01]$}: Newton tracks
        the informative branch when the basin is wide enough; $F_w$
        gets to $10^{-9}$ to $10^{-25}$ on some cells, but other
        cells slip into the trivial basin or land at intermediate
        non-FPs (high $F_\infty$, broken sym).
  \item \textbf{High $\tau$ ($\ge 0.05$)}: Newton plateaus much like
        Picard, with $F_w \sim 10^{-4}$ to $10^{-2}$ and frequent
        symmetry breaks. The underlying strict-$h{=}0$ operator is
        non-smooth here (Finding 9/17 in the main summary), and
        Newton's finite-difference Jacobian is noisy.
  \item \textbf{Cost}: Newton is $50\!-\!100\times$ more $F$-evals
        per cell than Picard (200--1000 vs $\sim 40$). Most of that
        is GMRES inner iterations.
\end{itemize}

\section{Comparison to Picard ladder}

\begin{center}\small\begin{tabular}{rrrrrr}\toprule
$\tau$ & Picard $F_w$ & Newton $F_\infty$ & Newton $F_w$ & Picard wall & Newton wall \\
\midrule
0.001  & 2.6e-15 & 3.3e-16 & 1.5e-67 & 4.8\,s & 0.3\,s \\
0.002  & 1.4e-14 & 3.3e-16 & 1.0e-54 & 0.2\,s & 0.1\,s \\
0.003  & 1.8e-2  & 2.8e-1  & 2.7e-10 & 0.3\,s & 3.1\,s \\
0.005  & 6.5e-4  & 1.5e-1  & 6.9e-10 & 0.3\,s & 4.5\,s \\
0.01   & 1.7e-4  & 1.4e-1  & 5.0e-5  & 0.2\,s & 3.0\,s \\
0.05   & 1.5e-2  & --      & --      & 0.2\,s & --     \\
1.0    & 1.5e-1  & --      & --      & 0.2\,s & --     \\
\bottomrule\end{tabular}\end{center}

(Full per-cell numbers in \texttt{logodds\_newton/ladder.json};
Newton entries for $\tau \ge 0.02$ have similar weighted residual but
the raw $F_\infty$ is dominated by symmetry-breaking artifacts in the
non-smooth strict-$h{=}0$ regime.)

\section{What Newton-IFT did and didn't deliver}

\textbf{Delivered.}
\begin{itemize}\itemsep1pt
  \item Confirmed that $P = 0.5$ is an exact FP (machine $\varepsilon$)
        across the whole $\tau$ ladder.
  \item Picked up the informative branch with weighted residual
        $F_w$ as low as $10^{-25}$ when the Picard warm-start was
        already close (e.g.\ $\tau = 0.00275$).
\end{itemize}

\textbf{Did NOT deliver.}
\begin{itemize}\itemsep1pt
  \item Machine-eps FP on the informative branch at high $\tau$.
        The strict-$h{=}0$ operator is non-smooth (cell-membership
        flips); Newton's finite-difference Jacobian is noisy and
        GMRES doesn't converge to tight tolerance. Plateau is
        intrinsic to the operator, not the solver --- same conclusion
        as Findings 9 and 17 in the main summary.
  \item Cost-effective replacement for Picard: Newton is
        $50\!-\!100\times$ slower per cell. Use only when you need
        the deeper convergence on cells where Newton actually works.
\end{itemize}

\textbf{Recommendation.} The Newton-IFT pass is most useful as a
\emph{refinement} step on top of Picard --- run Picard to navigate
basins and get close, then Newton for cells where the basin is
locally smooth enough for tight convergence. For the high-$\tau$
non-smooth regime, both methods hit the same operator wall;
breaking through requires fixing the operator (Finding 9 attempts;
smoothing introduces bias).

\section{Reproducibility}

\begin{itemize}\itemsep1pt
  \item Solver: \texttt{projects/REZN/solver\_code/highprec\_dd\_qd/dd\_k3\_logodds\_newton.py}.
  \item Per-cell FPs: \texttt{projects/REZN/solved\_fixed\_points/dd\_k3\_overnight/logodds\_newton/fp\_t*.npz}.
  \item Summary JSON: \texttt{ladder.json}.
  \item Live log: \texttt{/tmp/newton2.log} (per-cell printout).
\end{itemize}

\end{document}
